%% GENERATED by python/paper/build.py from results/ -- do not edit by hand.
%% Rebuild:  .venv\Scripts\python.exe python\paper\build.py --only OECDdata
%% Source:   data/oecdFigure1.csv, data/oecd26_theta.csv, data/oecd29_distribution.csv
\textit{Note:} 21 OECD members before 2000: those with informal employment below 25\% of employment (left out: Korea, Mexico, Turkey) where one public pay-as-you-go system pays the bulk of retirement benefits, that is, funded tiers pay less than a quarter of mandatory benefits in 2020 and the public pension is not tested against them (left out: Australia, Denmark, Iceland, Netherlands, Switzerland). The Bismarckian index is $I = 1/(2-\theta)$, with $\theta$ the earnings-linked share of the least-squares line through the public pension at 0.5, 1 and 2 times average earnings ($c/(a+c)$ of $p(y) = a + c\,y$), times one minus the share of public spending on minimum top-ups and tested benefits (Pensions at a Glance 2021). Pension spending is public cash spending on old age and survivors, 2020, in \% of GDP; population growth is the population of 2020 over that of 1990; high over low income is $y_3/y_1$, the relative income of the top of three income groups over that of the bottom one, the groups at the U.S.\ shares of the population (58\%, 18\%, 24\%, from the bottom), in the gross earnings of men working full time (OECD), 2019 (Poland 2018), as in the model.
